% !TeX root = ../informe.tex
% ===========================================================================
%  6. Reglas de negocio
%
%  Lo que se evalúa es la trazabilidad: cada regla, dónde se implementa y
%  cómo se comprueba.
%
%  Los \label{rn:NN} de la tabla son los que hacen funcionar \rn{NN} desde
%  el resto del documento: si añades filas, añade también su \label.
% ===========================================================================

\section{Reglas de negocio}
\label{sec:reglas-negocio}

\subsection{Catálogo de reglas}
\label{sec:catalogo-reglas}

\instruccion{Una fila por regla, con la columna de implementación rellenada
de verdad: «\id{ms-pedidos}, \id{PedidoService}» dice algo; «se valida en el
backend» no. Esta tabla es la que se mira para puntuar este criterio.}

\begin{table}[H]
  \centering
  \caption{Reglas de negocio y dónde se implementa cada una.}
  \label{tab:reglas}
  \begin{tabular}{C{1.6cm} L{7.2cm} L{4.4cm}}
    \toprule
    \textbf{ID} & \textbf{Descripción} & \textbf{Implementación} \\
    \midrule
    \textbf{RN-01}\label{rn:01} & \ph{Enunciado de la regla} & \ph{...} \\
    \textbf{RN-02}\label{rn:02} & \ph{Enunciado de la regla} & \ph{...} \\
    \textbf{RN-03}\label{rn:03} & \ph{Enunciado de la regla} & \ph{...} \\
    \bottomrule
  \end{tabular}
\end{table}

\subsection{Regla crítica}
\label{sec:regla-critica}

\instruccion{Elige la regla que más riesgo concentra — típicamente una que
falla solo bajo concurrencia o con datos límite — y dedícale espacio: qué
pasa si no se cumple, dónde se resuelve y por qué ahí. Renombra el apartado
con el nombre de la regla.}

\ph{Mecanismo elegido para \rn{01} y por qué resiste el caso adverso.}

\subsection{Máquina de estados}
\label{sec:estados}

\instruccion{Si alguna entidad del sistema pasa por estados, este es su
sitio. Un diagrama de estados pequeño comunica esto mejor que un párrafo; si
no hay estados que modelar, borra el apartado en vez de dejarlo vacío.}

\ph{Estados y transiciones permitidas.}

\subsection{Traducción de reglas a respuestas de la API}
\label{sec:errores}

\instruccion{Qué código devuelve la API cuando se viola cada regla. Importa
más de lo que parece: un 500 genérico ante una violación previsible es un
defecto, y un código específico con un mensaje claro es lo que permite que la
interfaz reaccione bien.}

\begin{table}[H]
  \centering
  \caption{Correspondencia entre violación de regla y respuesta HTTP.}
  \label{tab:errores}
  \begin{tabular}{C{1.6cm} C{2.4cm} L{8.6cm}}
    \toprule
    \textbf{Regla} & \textbf{HTTP} & \textbf{Código de error y significado} \\
    \midrule
    \ph{...} & \ph{...} & \ph{...} \\
    \bottomrule
  \end{tabular}
\end{table}
